\documentclass[11pt]{article}
\usepackage{amsmath,amssymb,amsthm,hyperref}
\newtheorem{theorem}{Theorem}
\newtheorem{proposition}[theorem]{Proposition}
\title{Weighted Affine Root Completions and Their First Variations}
\author{}
\date{November 23, 2026}
\begin{document}
\maketitle

\paragraph{Scope.}
The constructions below supply exact cancellation identities, with all
factors and all heads counted. They do not assert a universal small-head
representation. The missing universal choice of root lists and
perturbations remains a theorem-strength gap. The numerical search
described at the end tests a genuinely larger family of second lists
than coordinate permutations of the unweighted construction.

\begin{theorem}[Weighted interval completion]
Let $c_1,\ldots,c_n$ be positive integers, set
$s_i=c_1+\cdots+c_i$, $s_0=0$, and $M=s_n$.
Define $M+1$ affine forms by
\[
 u_0(x)=\sum_{i=1}^n c_i x_i,
 \qquad
 u_{i,j}(x)=j+(s_{i-1}+1)(1-x_i)+\sum_{k=i+1}^n c_kx_k,
 \quad 0\le j<c_i.
\]
For every $x\in\{0,1\}^n$, these values are a permutation of
$0,1,\ldots,M$. Consequently, as a polynomial in $s$,
\[
 (s+u_0(x))\prod_{i=1}^n\prod_{j=0}^{c_i-1}(s+u_{i,j}(x))
                         =\prod_{j=0}^{M}(s+j).
\tag{1}
\]
\end{theorem}
\begin{proof}
Induct on the number of coordinates. For no coordinates the list is
the singleton $0$. On appending coordinate $i$, the old list, which
has values $0,\ldots,s_{i-1}$, is translated by $c_ix_i$.
The new list consists of
$j+(s_{i-1}+1)(1-x_i)$ for $0\le j<c_i$.
If $x_i=0$, these two disjoint intervals are
$[0,s_{i-1}]$ and $[s_{i-1}+1,s_i]$.
If $x_i=1$, they are $[c_i,s_i]$ and $[0,c_i-1]$.
In both cases their union is $[0,s_i]$, with no repetition.
Subsequent steps translate every old entry by the same amount;
the resulting formulas are precisely the ones displayed.
Taking products proves (1).
\end{proof}

This includes the unweighted triangular list when every $c_i=1$.
For integer affine forms with nonzero coefficients of either sign,
complement each coordinate with a negative coefficient and subtract
the minimum value. The normalized form is then $u_0$ above, with
$c_i$ the absolute coefficient. Undoing the complements gives a
completion by affine forms. Coordinates with coefficient zero can be
omitted. The number of forms is one plus the sum of absolute
coefficients, not one plus the number of variables.

\begin{theorem}[An explicit first-variation compiler]
Let $u_0,\ldots,u_{m-1}$ and $v_0,\ldots,v_{m-1}$ be real affine
forms whose values are permutations of $0,\ldots,m-1$ at every
Boolean vertex. Let $a_h,b_h$ be real affine forms. Suppose the first
variation
\[
 R_x(s)=\left.\frac{\partial}{\partial\epsilon}\right|_{\epsilon=0}
 \left(\prod_h(s+u_h(x)+\epsilon a_h(x))
       -\prod_h(s+v_h(x)+\epsilon b_h(x))\right)
\]
is independent of $s$ and has strict prescribed signs $y(x)$.
Then these signs have a representation by $2m$ admissible heads with
constant numerators. If all supplied coefficients are rational, the
construction below chooses rational head parameters by finite exact
operations.
\end{theorem}
\begin{proof}
Put $P(s)=\prod_{j=0}^{m-1}(s+j)$ and
$C(x)=1+m\sum_i x_i$. Every constant coefficient of $u_h,v_h$
lies in $[0,m-1]$ and every variable coefficient in
$[-(m-1),m-1]$, by comparing the zero vertex to a unit vertex.
Thus every coefficient of $C+u_h,C+v_h$ is at least one.

Write $F_\epsilon(s,x)$ and $G_\epsilon(s,x)$ for the two products
in the hypothesis. A prime denotes differentiation with respect
to $s$, with $x$ held fixed. Define
\[
 Q_x(\epsilon)=
 (F_\epsilon G'_\epsilon-F'_\epsilon G_\epsilon)(C(x),x).
\]
At $\epsilon=0$ the two factors are $P$, so $Q_x(0)=0$.
Since $R_x$ is independent of $s$, product differentiation gives
\[
 Q'_x(0)=R_x P'(C(x)).
\tag{2}
\]
Here the prime on $Q_x$ means differentiation in $\epsilon$.
The number $C(x)$ is positive, and so is $P'(C(x))$.
In particular (2) has strict sign $y(x)$.

For complete quantitative choices, let $L\ge1$ bound the absolute
values of every coefficient of $a_h,b_h$. Expand
$Q_x(\epsilon)=\sum_{k=1}^{2m-1}q_{x,k}\epsilon^k$,
allowing zero coefficients. Set
\[
 \mu=\min_x|q_{x,1}|>0,
 \qquad T=\max_x\sum_{k=2}^{2m-1}|q_{x,k}|,
 \qquad
 \epsilon=\min\left\{\frac12,\frac1{2L},
                             \frac{\mu}{4(T+1)}\right\}>0.
\]
Each coefficient of every perturbed denominator is at least
$1-\epsilon L\ge1/2$, ensuring strict admissibility.
Also
\[
 \left|\frac{Q_x(\epsilon)}\epsilon-q_{x,1}\right|
       \le\epsilon\sum_{k=2}^{2m-1}|q_{x,k}|
       \le\frac{\mu T}{4(T+1)}<\mu.
\]
Consequently $Q_x(\epsilon)$ has sign $y(x)$ at all vertices.
Finally, logarithmic differentiation and positivity show that
\[
 -\sum_h\frac1{C(x)+u_h(x)+\epsilon a_h(x)}
 +\sum_h\frac1{C(x)+v_h(x)+\epsilon b_h(x)}
   =\frac{Q_x(\epsilon)}
           {F_\epsilon(C(x),x)G_\epsilon(C(x),x)}
\]
has precisely those signs. This is the stated $2m$-head score.
Every minimum, maximum, polynomial expansion, and comparison above
is finite; rational input coefficients give rational outputs.
\end{proof}

\begin{proposition}[Root-value test for an ideal first variation]
At a fixed vertex, let $h_j$ and $k_j$ be the unique indices with
$u_{h_j}(x)=v_{k_j}(x)=j$, and put
$d_j=(-1)^j j!(m-1-j)!$. Then $R_x(s)$ is the constant $r(x)$
if and only if
\[
 d_j\big(a_{h_j}(x)-b_{k_j}(x)\big)=r(x)
                   \quad(0\le j<m).
\tag{3}
\]
\end{proposition}
\begin{proof}
The degree of $R_x$ is at most $m-1$. Evaluation at $s=-j$
annihilates every summand of the product derivative except the
one corresponding to the unique root $j$. Its surviving product is
$\prod_{\ell\ne j}(\ell-j)=d_j$. Thus (3) is exactly the
statement that $R_x-r(x)$ vanishes at $m$ distinct points.
The degree bound proves the equivalence.
\end{proof}

\paragraph{Search and verification.}
The script \texttt{search\_permutation\_tangent\_smt.py} fixes the
first list to a weighted completion and lets the second list have
arbitrary integer slopes, with its constant terms fixed to
$0,\ldots,m-1$. This fixes only the row order at the zero vertex.
The integrality of the slopes follows already from the prescribed
integer root values at zero and at the unit vertices. It therefore
places no further restriction on affine second lists with this root set.
All cube vertices must evaluate to permutations. The perturbations
are arbitrary rational affine forms, and (3) is imposed with strict
parity signs and margin at least one. The solver's conditions use
linear arithmetic and conditional expressions. A returned witness
is checked independently using exact rational arithmetic at every
root and every vertex. A timeout is not an obstruction, and an
unsatisfiability claim would concern only this specified ansatz.

\paragraph{An explicit non-affine output.}
For four coordinates, take the weighted first list with
$(c_1,c_2,c_3,c_4)=(1,1,1,2)$ and the second list
\[
 \begin{aligned}
 v_0&=4x_1+x_2,&v_1&=1+x_1-x_2+3x_4,\\
 v_2&=2-2x_1+3x_2,&v_3&=3-x_2+2x_3-2x_4,\\
 v_4&=4-3x_1-x_2+x_4,&v_5&=5-x_2-2x_3-2x_4.
 \end{aligned}
\]
The complete rational perturbation coefficients are recorded in
\texttt{diagnostics/permutation-tangent-n4-weighted.json}.
Checking (3) gives, with $t_i=1-2x_i$,
\[
 R(x)=2t_3(9t_1t_2t_4+6t_1+6t_2+8t_4).
\]
It has the strict signs of four-bit parity. Indeed,
\[
 \left(\prod_{i=1}^4t_i\right)R(x)
       =2(9+6t_2t_4+6t_1t_4+8t_1t_2)\ge2.
\]
If $t_1=-t_2$, the expression inside parentheses is one.
If $t_1=t_2$, it is $17+12t_1t_4\ge5$.
The arithmetic verifier checks the root permutations and all 96
equalities in (3), and the compiler produces a twelve-head score
checked independently on all sixteen vertices.
This example only shows that the earlier affine-output obstruction
for an unweighted list against constant roots does not extend to
this enlarged family. Parity already has a four-head construction;
the example is not a head-complexity improvement.

\paragraph{Universal gap.}
No theorem here constructs such lists for every Boolean function
with $m=O(2^n/n^2)$. Giving that missing assertion another name would
not be progress on the target. The compiler is used only to assess
actual new identities when they are produced, not as an assumed
existence principle.
\end{document}
